\SGASection{4}{Appendix}
\SGARefTarget{sga3:xxv:section:4}

We used the following proposition in the course of the proof.

\begin{sgastatement}{Proposition 4.1}
Let \(S\) be a locally noetherian scheme of dimension at most \(1\);
let \(G\) and \(H\) be two \(S\)-group schemes of finite type; and let
\(f\colon G\to H\) be a quasi-finite separated group homomorphism.
If \(G\) is flat over \(S\),\footnote{Editors' note: We have added the
flatness hypothesis, which had been omitted.} then \(f\) is affine.
\end{sgastatement}
\SGARefTarget{sga3:xxv:proposition:4:1}

We shall give the proof only when \(G\) is smooth over \(S\), which is
indeed the hypothesis satisfied in the application above.

\SGANextSectionTarget{sga3:xxv:subsection:4-2}\SGASubsection{4.2}{}
\SGARefTarget{sga3:xxv:paragraph:4:2}

By EGA II, 1.6.4, we may suppose that \(S\) is reduced. By the usual
techniques of passage to the limit,\footnote{Editors' note: See
EGA IV\textsubscript{3}, 8.10.5(viii).} we may suppose \(S\) local. If
\(\dim(S)=0\), the assertion is trivial.\footnote{Editors' note:
Indeed, if \(S=\operatorname{Spec}(k)\), with \(k\) a field, then \(f\)
is the composite of the projection \(p\colon G\to G/\ker(f)\) and a
closed immersion \(i\). Since \(\ker(f)\) is finite over \(k\), \(p\)
is finite (Exposé VI\textsubscript{B}, 9.2).}
Suppose therefore that \(\dim(S)=1\). By faithfully flat descent, we may
suppose that \(S\) is complete with algebraically closed residue field.
After replacing \(S\) by its normalization \(\widetilde S\), if
necessary, we may suppose \(S\) normal (EGA II, 6.7.1, and
EGA \(0_{\mathrm{IV}}\), 23.1.5).\footnote{Editors' note:
Indeed, the irreducible components \(S_1,\ldots,S_r\) of \(S\) are
complete integral noetherian local schemes. By a theorem of Nagata
(EGA \(0_{\mathrm{IV}}\), 23.1.5), the normalization
\(\widetilde S_i\) is finite over \(S_i\), and hence
\(\widetilde S\) is finite over \(S\). A theorem of Chevalley
(EGA II, 6.7.1) then shows that if
\(f\times_S\widetilde S\) is affine, so is \(f\).}
We are thus reduced to the case in which \(S\) is the spectrum of a
complete discrete valuation ring \(A\) with algebraically closed
residue field.

\SGANextSectionTarget{sga3:xxv:subsection:4-3}\SGASubsection{4.3}{}
\SGARefTarget{sga3:xxv:paragraph:4:3}

Let \(\eta\) (respectively \(s\)) be the generic (respectively closed)
point of \(S\). Consider the image \(f_\eta(G_\eta)\) of \(G_\eta\)
in \(H_\eta\). It is a closed subgroup scheme of \(H_\eta\). Let
\(H'\) be the schematic closure of \(f_\eta(G_\eta)\) in \(H\).
Since \(H'\to H\) is affine (it is a closed immersion), we may replace
\(H\) by \(H'\), and thus suppose that \(H\) is flat over \(S\) and
that \(f_\eta\) is surjective. Since \(f_\eta\) is finite and \(G,H\)
are flat over \(S\), we have
\[
 \dim(G_s)=\dim(G_\eta)=\dim(H_\eta)=\dim(H_s).
\]

\SGANextSectionTarget{sga3:xxv:subsection:4-4}\SGASubsection{4.4}{}
\SGARefTarget{sga3:xxv:paragraph:4:4}

Let \(H_s^0,\ldots,H_s^n\) be the irreducible components of \(H_s\),
where \(H_s^0\) denotes the identity component, and let
\(z_0,\ldots,z_n\) be their generic points. Since each local ring
\(\mathcal O_{H,z_i}\) has dimension at most \(1\), the morphism
\[
 G\times_H\mathcal O_{H,z_i}\longrightarrow\mathcal O_{H,z_i}
\]
is affine, being quasi-finite and separated (see \SGARefLink{sga3:xvi:lemma:4:2}{Exposé XVI, Lemma 4.2}).
Thus \(G\) is affine over \(H\) in a neighborhood of \(z_i\). Let \(V\)
be the largest open subset of \(H\) over which \(G|_V\) is affine.
Then \(V\) contains every \(z_i\),\footnote{Editors' note: The following
sentence has been added.} and therefore contains at least one closed
point \(y_i\in H_s^i\). Moreover,
\(\kappa(y_i)=\kappa(s)\), since \(\kappa(s)\) is algebraically closed.

On the other hand, \(V\) is plainly stable under translation by any
element \(g\in G(S)\). We have \(\dim(G_s)=\dim(H_s)\), and \(f_s\) is
finite, so
\[
 f(s)\colon G_s^0(s)\longrightarrow H_s^0(s)
\]
is surjective. Since \(A\) is complete and \(G\) is smooth over \(S\),
the canonical map \(G^0(S)\to G_s^0(s)\) is surjective. Since
\(H_s^0(s)\) acts transitively on each \(H_s^i(s)\), it follows that
\(V\supset H_s(s)\), and therefore, since \(\kappa(s)\) is algebraically
closed, \(V\supset H_s\).\footnote{Editors' note: When \(G\) is not
assumed smooth over \(S\), one may argue as follows. Let \(h_0\) be a
rational point of \(H_s^0\), and let \(g_0\) be a rational point of
\(G_s^0\) such that \(f(g_0)=h_0\). By Exposé VI\textsubscript{B},
5.6.1, there is a commutative diagram
\[
\begin{tikzcd}[ampersand replacement=\&,column sep=4.0em,row sep=3.0em]
S'' \arrow[r,"g"] \arrow[d,"\pi"'] \arrow[dr,"\phi"]
  \& G \arrow[d]\\
S' \arrow[r,"w"] \& S
\end{tikzcd}
\]
in which \(w\) is étale and surjective, \(\pi\) is finite and
surjective, and \(\phi^{-1}(s)\) consists of a single point \(s''\)
such that \(g(s'')=g_0\). Then \(G_{S''}\to H_{S''}\) is affine above
a neighborhood of \(h_0y_i\), and the same is true of
\(G_{S'}\to H_{S'}\) (EGA II, 6.7.1), and hence of \(G\to H\) by
faithfully flat descent (EGA IV\textsubscript{2}, 2.7.1(xiii)).
Thus \(h_0y_i\in V\), from which it follows that \(V\) contains
\(H_s(s)\), and therefore \(H_s\).}

Clearly \(V\supset H_\eta\), since \(f_\eta\) is finite. Consequently,
\(V=H\). \SGAQED

\section*{Bibliography\footnote{Editors' note: Additional references
cited in this Exposé.}}
\addcontentsline{toc}{section}{Bibliography}
\SGARefTarget{sga3:xxv:bibliography}

\begin{description}[style=multiline,leftmargin=4.8em,labelwidth=4.0em]
\item[\textup{[BLie]}]
N.~Bourbaki, \textit{Groupes et algèbres de Lie}, Chapters I and
VII--VIII, Hermann, 1960 and 1975.

\item[\textup{[BT84]}]
F.~Bruhat and J.~Tits, \textit{Groupes réductifs sur un corps local II.
Schémas en groupes. Existence d'une donnée radicielle valuée},
Publ.\ Math.\ I.H.É.S. \textbf{60} (1984), 5--184.

\item[\textup{[DG70]}]
M.~Demazure and P.~Gabriel, \textit{Groupes algébriques},
Masson \& North-Holland, 1970.

\item[\textup{[Ko66]}]
B.~Kostant, \textit{Groups over \(\mathbb Z\)}, pp.~90--98 in
\textit{Algebraic groups and their discontinuous subgroups}
(A.~Borel and G.~D.~Mostow, eds.), Proc.\ Symp.\ Pure Math.\
\textbf{IX}, American Mathematical Society, 1966.

\item[\textup{[Lu09]}]
G.~Lusztig, \textit{Study of a \(\mathbb Z\)-form of the coordinate
ring of a reductive group}, J.\ Amer.\ Math.\ Soc.\ \textbf{22}
(2009), no.~3, 739--769.

\item[\textup{[Se66]}]
J.-P.~Serre, \textit{Algèbres de Lie semi-simples complexes},
Benjamin, 1966.

\item[\textup{[Sp98]}]
T.~A.~Springer, \textit{Linear algebraic groups}, 2nd ed.,
Birkhäuser, 1998.

\item[\textup{[Ti66]}]
J.~Tits, \textit{Sur les constantes de structure et le théorème
d'existence des algèbres de Lie semi-simples},
Publ.\ Math.\ I.H.É.S. \textbf{31} (1966), 21--58.
\end{description}
